% =====================================================================
%  preamble.tex — wiederverwendbare Praeambel fuer Formelsammlungen
%  Kurs: Elektronik / Bauelemente (etit 105, CAU Kiel)
%  Zweck: Abschreibvorlage. Alles muss von Hand reproduzierbar sein.
%  Farbcode (bewusst nur 3 Farben, mehr ist handschriftlich unrealistisch):
%     schwarz = Formel / Text
%     rot     = Ergebnis und Warnung
%     gruen   = Erkennungsmerkmal
% =====================================================================

% ---------- Sprache & Zeichensatz ----------
\usepackage[T1]{fontenc}
\usepackage[utf8]{inputenc}
\usepackage[ngerman]{babel}
\usepackage{lmodern}

% ---------- Mathematik ----------
\usepackage{amsmath,amssymb}
\usepackage{siunitx}
\sisetup{
  locale            = DE,
  per-mode          = symbol,
  range-units       = single,
  exponent-product  = \cdot,
  detect-weight     = true,
  detect-family     = true
}

% ---------- Layout ----------
\usepackage[
  a4paper,
  top=8mm, bottom=8mm, left=8mm, right=8mm,
  footskip=4mm, headsep=0mm
]{geometry}
\usepackage{multicol}
\setlength{\columnsep}{5mm}
\setlength{\columnseprule}{0.2pt}
\usepackage{microtype}
\usepackage{enumitem}
\setlist{nosep, leftmargin=3.6mm, labelsep=1.1mm, topsep=0.4mm}
\setlist[enumerate]{label=\textbf{\arabic*.}}
\setlist[itemize]{label=\raisebox{0.15ex}{\tiny$\blacksquare$}}
\usepackage{array}
\usepackage{booktabs}
\usepackage{tabularx}
\usepackage{ragged2e}

% ---------- Grafik ----------
\usepackage{tikz}
\usepackage[siunitx, RPvoltages]{circuitikz}
\usetikzlibrary{arrows.meta, calc, patterns, decorations.pathmorphing, positioning}
\ctikzset{
  bipoles/length      = 6mm,
  american,
  font                = \scriptsize,
  label/align        = straight
}
\usepackage{pgfplots}
\pgfplotsset{compat=1.18}

% ---------- Boxen ----------
\usepackage[most]{tcolorbox}
\tcbuselibrary{skins, breakable}

% ---------- Farben ----------
\usepackage{xcolor}
\definecolor{FSrot}{RGB}{200,25,25}      % Ergebnis / Warnung
\definecolor{FSgruen}{RGB}{0,120,55}     % Erkennungsmerkmal
\definecolor{FSgrau}{RGB}{120,120,120}   % Hilfslinien, Rahmen
\definecolor{FShell}{RGB}{247,247,247}   % dezenter Boxhintergrund

% Semantische Kurzbefehle.
% Alle drei funktionieren in Text- UND Mathemodus (\ifmmode-Weiche),
% damit man sie nicht danach unterscheiden muss, wo man gerade steht.
\newcommand{\erg}[1]{\ifmmode{\color{FSrot}#1}\else\textcolor{FSrot}{\textbf{#1}}\fi}
\newcommand{\warn}[1]{\ifmmode{\color{FSrot}#1}\else\textcolor{FSrot}{#1}\fi}
\newcommand{\ekm}[1]{\ifmmode{\color{FSgruen}#1}\else\textcolor{FSgruen}{#1}\fi}

% ---------- Schriftgroesse global ----------
% Sehr dichte Seite: 7pt Grundschrift, 8.2pt Zeilenabstand.
\renewcommand{\baselinestretch}{1.0}

% Schmale Spalten bei 7pt erzeugen sonst staendig Over-/Underfull-Boxen.
% \sloppy + emergencystretch loest das, ohne dass man Text kuerzen muss.
\sloppy
\setlength{\emergencystretch}{3em}
\hbadness=10000
\hfuzz=2pt

% ---------- Zaehler / Ueberschriften ohne Nummerierung ----------
\setcounter{secnumdepth}{-1}
\pagestyle{empty}

% =====================================================================
%  RANDNOTATION — Kernregel des Blattes
%  Jede algebraische Umformung bekommt einen Trennstrich + Notation.
%  Benutzung:  a = b \notat{\times R_E}
% =====================================================================
\newcommand{\notat}[1]{%
  \;\bigm|\;\text{\footnotesize\itshape #1}%
}
% Variante fuer abgesetzte Zeilen in align/aligned
\newcommand{\notatr}[1]{%
  &&\bigm|\;\text{\footnotesize\itshape #1}%
}

% =====================================================================
%  TASCHENRECHNER-HINWEIS
%  Erscheint als kleiner eingerueckter Block mit Rechnersymbol.
% =====================================================================
\newcommand{\tr}[1]{%
  \par\vspace{0.5mm}%
  \noindent\hspace*{1mm}%
  \begin{minipage}{\dimexpr\linewidth-2mm\relax}
    \footnotesize\itshape
    \textcolor{FSgrau}{\rule[-0.3mm]{0.25mm}{2.6mm}}\;%
    \textbf{TR:}\ #1
  \end{minipage}%
  \par\vspace{0.5mm}%
}

% =====================================================================
%  TCOLORBOX-STILE
% =====================================================================

% --- Basisstil, von allen anderen geerbt -----------------------------
\tcbset{
  FSbase/.style={
    boxrule      = 0.35mm,
    arc          = 0.6mm,
    left         = 1.1mm,
    right        = 1.1mm,
    top          = 0.9mm,
    bottom       = 0.9mm,
    boxsep       = 0mm,
    before skip  = 1.2mm,
    after skip   = 1.2mm,
    fonttitle    = \bfseries\footnotesize,
    fontupper    = \fontsize{7}{8.2}\selectfont,
    colback      = white,
    colframe     = black,
    coltitle     = white,
    enhanced,
    breakable    = false,
    % parbox=true ist noetig: tabularx misst innerhalb von multicols sonst
    % falsch und erzeugt eine Overfull-Box von 15pt (verifiziert).
    parbox       = true
  }
}

% --- Prio A: vollstaendiges Rezept -----------------------------------
% Kraeftiger Rahmen, Titelbalken schwarz. Handschriftlich = Doppelrahmen.
\newtcolorbox{RezeptA}[2][]{
  FSbase,
  colframe   = black,
  colbacktitle = black,
  boxrule    = 0.5mm,
  title      = {#2},
  #1
}

% --- Prio B: Formel + Erkennungsmerkmal, kein Zahlenbeispiel ---------
% Dezenter: duenner Rahmen, Titel schwarz auf hellgrau.
\newtcolorbox{RezeptB}[2][]{
  FSbase,
  colframe     = FSgrau,
  colbacktitle = FShell,
  coltitle     = black,
  boxrule      = 0.25mm,
  title        = {#2},
  #1
}

% --- Prio C: Einzeiler, nur Linie oben ------------------------------
\newtcolorbox{RezeptC}[2][]{
  FSbase,
  colframe     = FSgrau,
  colbacktitle = white,
  coltitle     = black,
  boxrule      = 0mm,
  frame hidden,
  borderline north = {0.25mm}{0mm}{FSgrau},
  title        = {#2},
  #1
}

% --- Neutraler Block ohne Prioritaetsanmutung (Konstanten, Tabellen) --
\newtcolorbox{FSblock}[2][]{
  FSbase,
  colframe     = black,
  colbacktitle = black,
  boxrule      = 0.35mm,
  title        = {#2},
  #1
}

% =====================================================================
%  STRUKTURBEFEHLE INNERHALB EINES REZEPTS
%  Reihenfolge laut Vorgabe:
%  1 Gesucht  2 Erkennungsmerkmal  3 Schrittkette  4 Zahlenbeispiel
%  5 TR-Hinweis  6 Variante
% =====================================================================
\newcommand{\Gesucht}[1]{%
  \par\noindent\textbf{Gesucht:}\ #1\par\vspace{0.3mm}}

\newcommand{\Erkennung}[1]{%
  \par\noindent\ekm{\textbf{Erkennungsmerkmal:}\ #1}\par\vspace{0.3mm}}

\newcommand{\Kette}{%
  \par\vspace{0.3mm}\noindent\textbf{Schrittkette}\par\vspace{0.2mm}}

\newcommand{\Beispiel}[1][]{%
  \par\vspace{0.6mm}\noindent\textbf{Zahlenbeispiel}%
  \ifx\relax#1\relax\else\ \textnormal{(#1)}\fi\par\vspace{0.2mm}}

\newcommand{\Variante}[1]{%
  \par\vspace{0.6mm}\noindent\textbf{Variante:}\ #1\par}

% Kleine Trennlinie innerhalb einer Box
\newcommand{\FSline}{\par\vspace{0.6mm}\noindent\textcolor{FSgrau}{\rule{\linewidth}{0.15mm}}\par\vspace{0.6mm}}

% Seitenkopf
\newcommand{\FSseite}[2]{%
  \noindent
  \begin{tcolorbox}[FSbase, colback=black, colframe=black,
                    left=1.5mm, right=1.5mm, top=0.5mm, bottom=0.5mm,
                    before skip=0mm, after skip=1.5mm]
    \color{white}\bfseries\small #1 \hfill \normalfont\footnotesize\itshape #2
  \end{tcolorbox}%
}

% =====================================================================
%  MATHE-ABKUERZUNGEN (Notation des Skripts)
% =====================================================================
\newcommand{\UT}{U_{T}}
\newcommand{\Uth}{U_{\mathrm{th}}}
\newcommand{\Cox}{C'_{\mathrm{ox}}}
\newcommand{\eps}{\varepsilon}
\newcommand{\dd}{\mathrm{d}}
\newcommand{\ee}{\mathrm{e}}
\renewcommand{\vec}[1]{\mathbf{#1}}

% Kompakte Einheitendarstellung im Fliesstext
\newcommand{\uu}[1]{\,\mathrm{#1}}

% =====================================================================
%  CIRCUITIKZ-SNIPPETS
%  Wiederverwendbare Schaltbilder. Alle mit Zaehlpfeilen fuer
%  Stroeme und Spannungen. Aufruf z.B. \snipEmitterverstaerker
%  Erweiterbar: neue Snippets hier unten anhaengen.
% =====================================================================

% --- Spannungspfeil-Helfer (offener Pfeil, handschriftlich machbar) ---
\newcommand{\Upfeil}[3]{% #1 von, #2 nach, #3 Label
  \draw[-{Latex[length=1.1mm]}, FSgrau] (#1) -- (#2) node[midway, right=0.4mm, black, font=\tiny] {#3};
}
